\documentclass[11pt]{article}

\usepackage[margin=1in]{geometry}
\usepackage[T1]{fontenc}
\usepackage{amsmath}
\usepackage{array}
\usepackage{amsfonts}
\usepackage{booktabs}
\usepackage{listings}

\newcolumntype{L}[1]{>{\raggedright\arraybackslash}p{#1}}
\lstset{
  basicstyle=\ttfamily\footnotesize,
  breaklines=true,
  columns=fullflexible,
  keepspaces=true
}

\title{Experiment 100 Human-Auditable Report\\
\large Square-Zero Nonunital Cofibrant Computation Through Weight 10}
\author{Jordan Quillen Homology Project}
\date{Run finished: 2026-06-03T04:47:17\\
Report updated with EXP101 generator-legality audit: 2026-06-03}

\begin{document}
\maketitle

\section*{Audit Scope}

This report summarizes the actual run of
\texttt{EXP-100-square-zero-nonunital-cofibrant-weight10}.  The target is the
nonunital Jordan algebra
\[
J=\mathrm{Jord}\langle x\rangle/(x^2)
\]
over \(\mathbb{Q}\).  The computation is truncated at product weight
\(w\leq 10\) and homological degree \(\leq 2\).  It computes the bounded cell
layers \(V_0,V_1,V_2\), does not apply the functor \(Q\), and does not
construct any layers \(V_3,V_4,\ldots\).

The value of the degree-2 layer was not assumed before the run.  In
particular, the plan and runner treated \(\widetilde{J}_2\) as unknown data to
be discovered by exact linear algebra.  The setup file
\texttt{expected.json} records only the experiment bounds and explicitly marks
computed cells, differentials, predicted raw cycle counts, predicted minimal
cell tables, and outcomes as fields not to prefill.

\section*{Primary Audit Artifacts}

The full machine-readable audit trail is kept in JSON.  This TeX file is a
human-readable report and intentionally does not duplicate every basis vector
or coordinate array.

\begin{lstlisting}
Plan:
researchplan/subplan100.md

Runner:
experiments/100-square-zero-nonunital-cofibrant-weight10/run.py

Setup-only expectation file:
experiments/100-square-zero-nonunital-cofibrant-weight10/expected.json

Primary result object:
experiments/100-square-zero-nonunital-cofibrant-weight10/results.json

Per-weight audit data:
experiments/100-square-zero-nonunital-cofibrant-weight10/data/by_weight_W10.json

Cell audit data:
experiments/100-square-zero-nonunital-cofibrant-weight10/data/cells_W10.json

Run log:
experiments/100-square-zero-nonunital-cofibrant-weight10/logs/run_W10.log

Post-run generator-legality audit:
experiments/101-one-generator-generator-legality-audit/results.json
\end{lstlisting}

The reproducing command recorded for the run is:
\begin{lstlisting}
python experiments/100-square-zero-nonunital-cofibrant-weight10/run.py --max-weight 10
\end{lstlisting}

Raw target field from \texttt{results.json}:
\begin{lstlisting}
Jord<x>/(x^2)
\end{lstlisting}

\section*{Backend and Conventions}

The run used exact rational arithmetic.  The backend convention recorded in
\texttt{results.json} is:
\begin{lstlisting}
ordinary commutative nonassociative Jordan algebra modulo the linearized Jordan identity, with a signed derivation differential
\end{lstlisting}

The matrix convention is
\texttt{source\_rows\_target\_coordinates}.  The output is therefore relative
to this ordinary backend convention.  A separate theory check is still needed
before identifying this computation with a fully graded operadic dg Jordan
convention.

\section*{Computation Procedure}

The computation proceeds weight by weight.

\begin{enumerate}
\item Build \(V_0\) from the presentation generator \(x\), mapped to
\(\bar{x}\) in \(J\).
\item For each weight \(w\leq 10\), compute the presentation kernel in the
ordinary nonunital Jordan backend.  A new \(V_1\) relation cell is accepted
only when its differential maps to zero in \(J\) and its representative is
independent modulo the prior closure generated by previously accepted
relations.
\item After the accepted \(V_1\)-layer is known up to the current weight, form
the old degree-one complex
\[
C^{old}_{2,w}\xrightarrow{d^{old}_2}C_{1,w}\xrightarrow{d_1}C_{0,w}.
\]
\item Compute \(Z_{1,w}=\ker d_1\), \(B^{old}_{1,w}=\mathrm{im}(d^{old}_2)\),
and \(H^{old}_{1,w}=Z_{1,w}/B^{old}_{1,w}\).  Candidate \(V_2\)
representatives are selected from certified cycles whose RREF remainder modulo
old boundaries is nonzero and independent from previously selected
representatives.
\item Record chain, rank-nullity, cycle, non-boundary, and independence checks
in \texttt{results.json} and \texttt{data/by\_weight\_W10.json}.
\end{enumerate}

\section*{Run Summary}

\begin{center}
\begin{tabular}{ll}
\toprule
Field & Value \\
\midrule
Status & \texttt{run} \\
Passed & \texttt{true} \\
Requested maximum weight & \(10\) \\
Completed weights & \(1,2,3,4,5,6,7,8,9,10\) \\
Skipped weights & none \\
Failed weights & none \\
Runtime recorded in \texttt{results.json} & \(0.116\) seconds \\
Total \(V_0\) cells & \(1\) \\
Total \(V_1\) cells & \(3\) \\
Total \(V_2\) cells & \(5\) \\
\bottomrule
\end{tabular}
\end{center}

\section*{\(V_1\) Discovery Audit}

The following table records the relation-discovery pass.  The column
\(\dim A_{0,w}\) is the source free-Jordan basis dimension in weight \(w\);
\(\dim\ker\epsilon_w\) is the presentation-kernel dimension in that weight;
``closure rank'' is the rank of the relation closure after accepting any new
minimal \(V_1\) cells at that weight.

Raw \(V_1\) cell names from \texttt{data/cells\_W10.json}:
\begin{lstlisting}
r1_00001_w2
r1_00002_w10
r1_00003_w10
\end{lstlisting}

\begin{center}
\begin{tabular}{rrrrL{0.27\linewidth}}
\toprule
Weight & \(\dim A_{0,w}\) & \(\dim\ker\epsilon_w\) & Closure rank & New \(V_1\) cells \\
\midrule
1 & 1 & 0 & 0 & none \\
2 & 1 & 1 & 1 & \texttt{r1\_00001\_w2} \\
3 & 1 & 1 & 1 & none \\
4 & 1 & 1 & 1 & none \\
5 & 1 & 1 & 1 & none \\
6 & 2 & 2 & 2 & none \\
7 & 2 & 2 & 2 & none \\
8 & 5 & 5 & 5 & none \\
9 & 8 & 8 & 8 & none \\
10 & 17 & 17 & 17 & \texttt{r1\_00002\_w10}, \texttt{r1\_00003\_w10} \\
\bottomrule
\end{tabular}
\end{center}

\paragraph{Why new \(V_1\) cells appear at weight \(10\).}
The two weight-10 relation cells are not an imposed prediction.  In weight
\(10\), the presentation kernel has dimension \(17\).  Before adding the two
weight-10 relation cells, the closure generated by earlier \(V_1\) relations
has rank \(15\).  Thus two independent kernel directions remain outside the
old closure.  The runner chooses two representatives with nonzero RREF
remainders, recorded below as sparse coordinates \(\{0:1\}\) and
\(\{8:1\}\), and only then the closure rank reaches \(17\).

\subsection*{Accepted \(V_1\) Cells}

\paragraph{\texttt{r1\_00001\_w2}, weight 2.}
Differential:
\begin{lstlisting}
(x*x)
\end{lstlisting}
Certificate summary: maps to \(0\) in \(J\); RREF remainder modulo prior
closure is nonzero; accepted as minimal in weight 2.

\paragraph{\texttt{r1\_00002\_w10}, weight 10.}
Differential:
\begin{lstlisting}
((((((x*x)*x)*x)*x)*x)*((x*x)*(x*x)))
\end{lstlisting}
Certificate summary: maps to \(0\) in \(J\); prior closure rank is \(15\);
RREF remainder modulo prior closure is nonzero with sparse coordinate
\(\{0:1\}\); accepted as minimal in weight 10.

\paragraph{\texttt{r1\_00003\_w10}, weight 10.}
Differential:
\begin{lstlisting}
(((((x*x)*x)*x)*x)*(((x*x)*x)*(x*x)))
\end{lstlisting}
Certificate summary: maps to \(0\) in \(J\); prior closure rank is \(16\);
RREF remainder modulo prior closure is nonzero with sparse coordinate
\(\{8:1\}\); accepted as minimal in weight 10.

\section*{Weight-by-Weight Homology Audit}

For each completed weight, the runner records \(C_{0,w}\), \(C_{1,w}\), the
old \(C^{old}_{2,w}\), the ranks of \(d_1\) and \(d^{old}_2\), and
\(\dim H^{old}_{1,w}\).  The check
\texttt{dim\_H1\_old\_matches\_V2\_count} is true in every row below.

\begin{center}
\small
\begin{tabular}{rrrrrrrr}
\toprule
\(w\) & \(\dim C_0\) & \(\dim C_1\) & \(\dim C^{old}_2\) &
\(\mathrm{rank}\,d_1\) & \(\dim Z_1\) & \(\mathrm{rank}\,d^{old}_2\) &
\(\dim H^{old}_1\) \\
\midrule
1 & 1 & 0 & 0 & 0 & 0 & 0 & 0 \\
2 & 1 & 1 & 0 & 1 & 0 & 0 & 0 \\
3 & 1 & 1 & 0 & 1 & 0 & 0 & 0 \\
4 & 1 & 2 & 1 & 1 & 1 & 0 & 1 \\
5 & 1 & 2 & 2 & 1 & 1 & 1 & 0 \\
6 & 2 & 4 & 4 & 2 & 2 & 2 & 0 \\
7 & 2 & 7 & 7 & 2 & 5 & 4 & 1 \\
8 & 5 & 14 & 17 & 5 & 9 & 8 & 1 \\
9 & 8 & 29 & 37 & 8 & 21 & 21 & 0 \\
10 & 17 & 65 & 90 & 17 & 48 & 46 & 2 \\
\bottomrule
\end{tabular}
\end{center}

\begin{center}
\begin{tabular}{rL{0.38\linewidth}L{0.38\linewidth}}
\toprule
\(w\) & New \(V_1\) & New \(V_2\) \\
\midrule
1 & none & none \\
2 & \texttt{r1\_00001\_w2} & none \\
3 & none & none \\
4 & none & \texttt{s2\_00001\_w4} \\
5 & none & none \\
6 & none & none \\
7 & none & \texttt{s2\_00002\_w7} \\
8 & none & \texttt{s2\_00003\_w8} \\
9 & none & none \\
10 & \texttt{r1\_00002\_w10}, \texttt{r1\_00003\_w10} &
\texttt{s2\_00004\_w10}, \texttt{s2\_00005\_w10} \\
\bottomrule
\end{tabular}
\end{center}

\section*{Candidate \(V_2\) Representatives Selected by the Old-Homology Audit}

Each listed \(V_2\) candidate below has three recorded EXP100 certificates:
the chosen differential is a cycle (\(d_1 z=0\)); its remainder modulo old
boundaries is nonzero; and the selected representative is independent modulo
old boundaries plus previously selected \(V_2\) representatives.  These
certificates are old-complex certificates only; they are not, by themselves,
strict dg attaching-legality certificates.

Raw \(V_2\) cell names from \texttt{data/cells\_W10.json}:
\begin{lstlisting}
s2_00001_w4
s2_00002_w7
s2_00003_w8
s2_00004_w10
s2_00005_w10
\end{lstlisting}

\paragraph{\texttt{s2\_00001\_w4}, weight 4.}
Differential:
\begin{lstlisting}
-((r1_00001_w2*x)*x) + ((x*x)*r1_00001_w2)
\end{lstlisting}
Cycle normal form: \texttt{0}.  Boundary remainder modulo \(B^{old}\):
\begin{lstlisting}
-((r1_00001_w2*x)*x) + ((x*x)*r1_00001_w2)
\end{lstlisting}
Certificate type: \texttt{rref\_remainder}; remainder nonzero; independent
modulo \(B^{old}\) plus prior selected representatives.

\paragraph{\texttt{s2\_00002\_w7}, weight 7.}
Differential:
\begin{lstlisting}
-(((((x*x)*x)*x)*x)*r1_00001_w2) + ((((x*x)*x)*x)*(r1_00001_w2*x))
\end{lstlisting}
Cycle normal form: \texttt{0}.  Boundary remainder modulo \(B^{old}\):
\begin{lstlisting}
-1/2*(((r1_00001_w2*x)*(x*x))*(x*x)) + 1/2*(((x*x)*r1_00001_w2)*((x*x)*x))
\end{lstlisting}
Certificate type: \texttt{rref\_remainder}; remainder nonzero; independent
modulo \(B^{old}\) plus prior selected representatives.

\paragraph{\texttt{s2\_00003\_w8}, weight 8.}
Differential:
\begin{lstlisting}
-((((((x*x)*x)*x)*x)*x)*r1_00001_w2) + 1/2*(((((x*x)*(x*x))*x)*x)*r1_00001_w2) - 1/2*(((((x*x)*x)*x)*r1_00001_w2)*(x*x)) + (((((x*x)*x)*x)*x)*(r1_00001_w2*x))
\end{lstlisting}
Cycle normal form: \texttt{0}.  Boundary remainder modulo \(B^{old}\):
\begin{lstlisting}
1/4*((((x*x)*x)*((x*x)*x))*r1_00001_w2) - 1/4*((((x*x)*x)*(r1_00001_w2*x))*(x*x))
\end{lstlisting}
Certificate type: \texttt{rref\_remainder}; remainder nonzero; independent
modulo \(B^{old}\) plus prior selected representatives.

\paragraph{\texttt{s2\_00004\_w10}, weight 10.}
Differential:
\begin{lstlisting}
2*((((((((x*x)*x)*x)*x)*x)*x)*x)*r1_00001_w2) - (((((((x*x)*x)*x)*x)*x)*r1_00001_w2)*(x*x)) - 2*(((((((x*x)*x)*x)*x)*x)*x)*(r1_00001_w2*x)) + r1_00002_w10
\end{lstlisting}
Cycle normal form: \texttt{0}.  Boundary remainder modulo \(B^{old}\):
\begin{lstlisting}
-(((((r1_00001_w2*x)*x)*x)*x)*((x*x)*(x*x))) + r1_00002_w10
\end{lstlisting}
Certificate type: \texttt{rref\_remainder}; remainder nonzero; independent
modulo \(B^{old}\) plus prior selected representatives.

\paragraph{\texttt{s2\_00005\_w10}, weight 10.}
Differential:
\begin{lstlisting}
((((((((x*x)*x)*x)*x)*x)*x)*x)*r1_00001_w2) + 1/12*(((((((x*x)*(x*x))*x)*x)*x)*x)*r1_00001_w2) - 1/12*(((((((x*x)*x)*(x*x))*x)*x)*x)*r1_00001_w2) + 1/4*(((((((x*x)*x)*x)*x)*x)*r1_00001_w2)*(x*x)) - 1/3*(((((((x*x)*x)*x)*x)*x)*x)*(r1_00001_w2*x)) - 1/3*((((((x*x)*(x*x))*x)*x)*x)*(r1_00001_w2*x)) - 11/12*((((((x*x)*x)*x)*x)*(x*x))*(r1_00001_w2*x)) - 2/3*((((((x*x)*x)*x)*x)*x)*((r1_00001_w2*x)*x)) + r1_00003_w10
\end{lstlisting}
Cycle normal form: \texttt{0}.  Boundary remainder modulo \(B^{old}\):
\begin{lstlisting}
-((((r1_00001_w2*x)*x)*x)*(((x*x)*x)*(x*x))) + r1_00003_w10
\end{lstlisting}
Certificate type: \texttt{rref\_remainder}; remainder nonzero; independent
modulo \(B^{old}\) plus prior selected representatives.

\section*{Global Checks Recorded by the Runner}

\begin{center}
\begin{tabular}{L{0.70\linewidth}c}
\toprule
Check & Result \\
\midrule
\texttt{all\_V2\_differentials\_are\_cycles} & true \\
\texttt{all\_V2\_representatives\_have\_nonzero\_B\_old\_remainder} & true \\
\texttt{all\_V2\_representatives\_independent\_mod\_B\_old} & true \\
\texttt{all\_completed\_weights\_passed} & true \\
\texttt{all\_relation\_cells\_map\_to\_kernel} & true \\
\texttt{applies\_Q\_is\_false} & true \\
\texttt{chain\_condition\_checks\_passed} & true \\
\texttt{constructs\_higher\_cells\_V3\_plus\_is\_false} & true \\
\texttt{exact\_rational\_arithmetic} & true \\
\texttt{expected\_setup\_has\_no\_prefilled\_results} & true \\
\texttt{rank\_nullity\_checks\_passed} & true \\
\texttt{report\_language\_has\_no\_forbidden\_global\_phrasing} & true \\
\bottomrule
\end{tabular}
\end{center}

\section*{Post-Run Generator-Legality Audit (EXP101)}

EXP101 was run after EXP100 to test whether the generators selected in EXP100
are legal as strict dg attaching data in the same bounded backend.  The audit
uses EXP100's \texttt{results.json} and \texttt{data/cells\_W10.json} as
inputs, extends the legality check through weight \(12\), does not apply \(Q\),
does not construct \(V_3,V_4,\ldots\), and makes no homology-killing claim.

\begin{center}
\begin{tabular}{L{0.64\linewidth}c}
\toprule
EXP101 check & Result \\
\midrule
Requested generator-legality bound & \(12\) \\
Reference EXP100 run passed & \texttt{true} \\
Declared generators satisfy \(d^2=0\) & \texttt{true} \\
\(V_1\) relation generators legal through weight \(12\) & \(3/3\) \\
\(V_2\) generator-level differentials are cycles & \(5/5\) \\
\(V_2\) candidates strictly attachable through weight \(12\) & \(0/5\) \\
Strict attachability failures recorded & \(5\) \\
\bottomrule
\end{tabular}
\end{center}

The first strict-attachability defects recorded by EXP101 are:

\begin{center}
\begin{tabular}{rll}
\toprule
Candidate & First defect weight & Multiplier \\
\midrule
\texttt{s2\_00001\_w4} & \(6\) & \texttt{(x*x)} \\
\texttt{s2\_00002\_w7} & \(8\) & \texttt{x} \\
\texttt{s2\_00003\_w8} & \(9\) & \texttt{x} \\
\texttt{s2\_00004\_w10} & \(11\) & \texttt{x} \\
\texttt{s2\_00005\_w10} & \(11\) & \texttt{x} \\
\bottomrule
\end{tabular}
\end{center}

Thus EXP101 confirms the EXP100 \(V_1\) relation generators through the tested
range, but it does not validate the five chosen \(V_2\) representatives as
strict attaching differentials.  A \(V_2\) obstruction here means that the
chosen representative fails the bounded multiplicative strictness test in
\(A^{(1)}\).  It does not rule out a different representative, a corrected
lift, or behavior outside the tested bound.

\section*{Result Interpretation}

Within the recorded truncation \(w\leq 10\), the run found:
\[
|V_0|=1,\qquad |V_1|=3,\qquad |V_2|=5.
\]
The five \(V_2\) cells occur in weights \(4,7,8,10,10\).  The two weight-10
\(V_2\) cells explicitly involve the two new weight-10 \(V_1\) relation cells
\texttt{r1\_00002\_w10} and \texttt{r1\_00003\_w10}.

After EXP101, the safe interpretation is sharper: the three \(V_1\) relation
generators are legal through the tested weight \(12\), while the five EXP100
\(V_2\) records should be read as old-complex cycle representatives rather than
as validated strict \(V_2\) attaching cells.  This is bounded, backend-relative
computational evidence.  It does not prove global stabilization, does not
assert absence of cells above weight 10, does not construct \(V_3\) or higher
layers, and does not identify the ordinary backend output with a fully graded
operadic dg Jordan model without further theory work.

\end{document}
